\section{2025-01-15}

\startbuffer[physik:atomhuellen:energienevau]
\startMPpage
% Einheit definieren
u := 0.5cm;

color purple;
purple := (1, 0, 1); % RGB-Werte für Lila

% Zeichne die Achse und Energieniveaus
drawarrow (0,0)--(0,10u) withpen pencircle scaled 0.7pt;
label.top(btex $E$ etex, (0,10.5u));

% Energieniveaus
draw (0,3u)--(4.5u,3u) withpen pencircle scaled 0.5pt;
draw (0,6u)--(4.5u,6u) withpen pencircle scaled 0.5pt;
draw (0,9u)--(4.5u,9u) withpen pencircle scaled 0.5pt;

label.lft(btex $E_1$ etex, (6u,3u));
label.lft(btex $E_2$ etex, (6u,6u));
label.lft(btex $E_3$ etex, (6u,9u));

% Absorptionspfeile
drawarrow (2u,3u)--(2u,9u) withcolor blue;
label.lft(btex $\lambda_1$ etex, (2u,4.5u));

drawarrow (3u,3u)--(3u,6u) withcolor purple;
label.rt(btex $\lambda_2$ etex, (3u,4.5u));

drawarrow (4u,6u)--(4u,9u) withcolor green;
label.rt(btex $\lambda_3$ etex, (4u,7.5u));

% Zeichne die Achse und Energieniveaus
drawarrow (10u,0)--(10u,10u) withpen pencircle scaled 0.7pt;
label.top(btex $E$ etex, (10u,10.5u));

% Energieniveaus
draw (10u,3u)--(14.5u,3u) withpen pencircle scaled 0.5pt;
draw (10u,6u)--(14.5u,6u) withpen pencircle scaled 0.5pt;
draw (10u,9u)--(14.5u,9u) withpen pencircle scaled 0.5pt;

label.lft(btex $E_1$ etex, (9u,3u));
label.lft(btex $E_2$ etex, (9u,6u));
label.lft(btex $E_3$ etex, (9u,9u));

% Emissionsteil
drawarrow (12u,9u)--(12u,3u) withcolor blue;
label.lft(btex $\lambda_1'$ etex, (12u,4.5u));

drawarrow (13u,9u)--(13u,6u) withcolor purple;
label.rt(btex $\lambda_2'$ etex, (13u,7.5u));

drawarrow (14u,6u)--(14u,3u) withcolor green;
label.rt(btex $\lambda_3'$ etex, (14u,4.5u));

% Beschriftung Absorption und Emission
label.top(btex Absorption etex, (2u,12u));
label.top(btex Emission etex, (12u,12u));

% Formeln Absorption
label.lft(btex $\Delta E = E_3 - E_1$ etex, (0,1.5u));
label.lft(btex $E_{\text{Photon}} = h \cdot f = \frac{h \cdot c}{\lambda_1}$ etex, (0,0.5u));

% Formeln Emission
label.rt(btex $\Delta E = E_3 - E_1$ etex, (10u,1.5u));
label.rt(btex $E_{\text{Photon}} = h \cdot f = \frac{h \cdot c}{\lambda_1'}$ etex, (10u,0.5u));

\stopMPpage
\stopbuffer

\placefigure[force][physik:atomhuellen:energienevau]{Name}{
\typesetbuffer[physik:atomhuellen:energienevau]
}

\startframedtask
vlg. Buch S. 228
\stopframedtask

\section{2025-01-15 Fluoreszens und Phosphoreszens}

\startframedtipp
Bei Fluoreszens gibt es das angeregte Teilchen eine geringere Energie ab, als möglich wäre, bei Fostphoreszens findet die Emission verzögert in kleineren Teilschritten statt.
\stopframedtipp


\startbuffer[physik:atomhuellen:fluoreszens]
\startMPpage
% Einheit definieren
u := 0.5cm;

color purple;
purple := (1, 0, 1); % RGB-Werte für Lila

% Zeichne die Achse und Energieniveaus
drawarrow (0,0)--(0,10u) withpen pencircle scaled 0.7pt;
label.top(btex $E$ etex, (0,10.5u));

% Energieniveaus
draw (0,3u)--(4.5u,3u) withpen pencircle scaled 0.5pt;
draw (0,6u)--(4.5u,6u) withpen pencircle scaled 0.5pt;
draw (0,7.5u)--(4.5u,7.5u) withpen pencircle scaled 0.5pt;
draw (0,8.25u)--(4.5u,8.25u) withpen pencircle scaled 0.5pt;
draw (0,9u)--(4.5u,9u) withpen pencircle scaled 0.5pt;

% label.lft(btex $E_1$ etex, (6u,3u));
% label.lft(btex $E_2$ etex, (6u,6u));
% label.lft(btex $E_3$ etex, (6u,9u));

% Absorptionspfeile
drawarrow (2u,3u)--(2u,9u) withcolor blue;
label.lft(btex $E_{ab}$ etex, (2u,4.5u));

drawarrow (3u,9u)--(3u,6u) withcolor purple;
label.rt(btex $E_{emi}$ etex, (3u,7.5u));

% drawarrow (4u,6u)--(4u,9u) withcolor green;
% label.rt(btex $\lambda_3$ etex, (4u,7.5u));

% Zeichne die Achse und Energieniveaus
drawarrow (10u,0)--(10u,10u) withpen pencircle scaled 0.7pt;
label.top(btex $E$ etex, (10u,10.5u));

% Energieniveaus
draw (10u,3u)--(14.5u,3u) withpen pencircle scaled 0.5pt;
draw (10u,6u)--(14.5u,6u) withpen pencircle scaled 0.5pt;
draw (10u,7.5u)--(14.5u,7.5u) withpen pencircle scaled 0.5pt;
draw (10u,8.25u)--(14.5u,8.25u) withpen pencircle scaled 0.5pt;
draw (10u,9u)--(14.5u,9u) withpen pencircle scaled 0.5pt;

% Emissionsteil
drawarrow (12u,3u)--(12u,8.25u) withcolor blue;
label.lft(btex $E_{ab}$ etex, (12u,4.5u));

drawarrow (13u,8.25u)--(13u,7.5u) withcolor purple;

drawarrow (14u,7.5u)--(14u,6u) withcolor purple;

drawarrow (15u,6u)--(15u,3u) withcolor purple;

label.rt(btex $E_{emi 1}$ etex, (16u,7.825u));
label.rt(btex $E_{emi 2}$ etex, (16u,6.75u));
label.rt(btex $E_{emi 3}$ etex, (16u,4.5u));


% Beschriftung Absorption und Emission
label.top(btex Absorption etex, (2u,12u));
label.top(btex Emission etex, (12u,12u));

% Formeln Absorption
label.lft(btex Fluoreszens etex, (0,1.5u));
label.lft(btex $\lambda_{\text{rein}} < \lambda_{\text{raus}}$ etex, (0,0.5u));

% Formeln Emission
label.rt(btex Phosphoreszens etex, (10u,1.5u));
label.rt(btex $\lambda_{\text{rein}} < \lambda_{\text{raus }3} < \lambda_2 < \lambda_1$ etex, (10u,0.5u));

\stopMPpage
\stopbuffer

\placefigure[force][physik:atomhuellen:fluoreszens]{Name}{
\typesetbuffer[physik:atomhuellen:fluoreszens]
}

\section{2025-01-17 Leuchtsoffröhren}


Inhalt z.B.:
\startitemize[packed]
\item Neonröhren
\item Quecksilber
\item Argon
\item Xenon
\stopitemize

Leuchtstoff wird mit bewegten Elektronen angeregt (vlg. Spektrallampe)

Menschen sind kontinuirliche Spektrum der Sonne gewöhnt, mit Linien ist das Sehen von Farben nicht sinnvoll möglich.

Um einem kontinuierlichen Spektrum nahezukommen, werden Leuchtstofflampen mit unterschiedlichen fluoreszierenden Stoffen beschichtet. 
